% ==============================================================================
% dxh_pilot.tex — DXH workshop Part 1: piloting experience
% DXH講座 ①: ドローン操縦体験
% ==============================================================================

% ------------------------------------------------------------------
\begin{frame}{① 操縦のしかた（スティックモード3）}
  \centering
  \begin{adjustbox}{max width=0.88\textwidth, max height=0.72\textheight, center}
    % ==============================================================================
% dxh_stick_mode3.tex — DXH workshop: stick assignment in Mode 3
% DXH講座: スティックモード3の操作割当図
%
% Axis mapping verified against controller firmware (main.cpp:1545):
%   Mode 3: Throttle=right-Y, Rudder(yaw)=right-X,
%           Elevator(pitch)=left-Y, Aileron(roll)=left-X
% 軸割当はコントローラファーム (main.cpp:1545) で裏取り済み:
%   Mode 3: スロットル=右Y, ラダー(ヨー)=右X, エレベータ(ピッチ)=左Y,
%           エルロン(ロール)=左X
% ==============================================================================
\documentclass[border=10pt]{standalone}
\usepackage{tikz}
\usetikzlibrary{arrows.meta, positioning}
\usepackage{luatexja}
\usepackage[match]{luatexja-fontspec}
% Cross-platform font: Hiragino (macOS) → Yu Gothic (Windows) → Noto (Linux)
\usepackage{ifthen}
\newboolean{jfontset}\setboolean{jfontset}{false}
\IfFontExistsTF{Hiragino Kaku Gothic ProN}{%
  \setmainjfont{Hiragino Kaku Gothic ProN}%
  \setsansjfont{Hiragino Kaku Gothic ProN}%
  \setboolean{jfontset}{true}}{}
\ifthenelse{\boolean{jfontset}}{}{%
  \IfFontExistsTF{Yu Gothic}{%
    \setmainjfont{Yu Gothic}%
    \setsansjfont{Yu Gothic}%
    \setboolean{jfontset}{true}}{}}
\ifthenelse{\boolean{jfontset}}{}{%
  \IfFontExistsTF{Noto Sans CJK JP}{%
    \setmainjfont{Noto Sans CJK JP}%
    \setsansjfont{Noto Sans CJK JP}}{}}

\begin{document}

% Colors defined after \begin{document} (standalone compatibility, rule T9)
% 色定義は \begin{document} の後（standalone 互換, ルール T9）
\definecolor{cstick}{RGB}{60, 60, 60}
\definecolor{cmove}{RGB}{0, 100, 180}
\definecolor{cvert}{RGB}{180, 60, 40}

\begin{tikzpicture}[
    stickname/.style={font=\bfseries\large, text=black!70, align=center},
    axislabel/.style={font=\normalsize, align=center},
    arr/.style={-{Stealth[length=3mm]}, line width=1.8pt, #1},
    note/.style={font=\small, text=black!55, align=center},
  ]

  % === Left stick: horizontal movement / 左スティック: 水平方向の移動 ===
  \node[stickname] at (-6.0, 4.1) {左スティック};
  \draw[cstick, line width=2pt, fill=black!25] (-6.0, 0) circle (1.4);
  \draw[cstick, line width=1pt, fill=black!40] (-6.0, 0) circle (0.9);

  \draw[arr=cmove] (-6.0, 1.6) -- (-6.0, 2.8)
    node[axislabel, text=cmove, above=0.5mm] {前進};
  \draw[arr=cmove] (-6.0, -1.6) -- (-6.0, -2.8)
    node[axislabel, text=cmove, below=0.5mm] {後退};
  \draw[arr=cmove] (-7.6, 0) -- (-8.8, 0)
    node[axislabel, text=cmove, left=0.5mm] {左へ移動};
  \draw[arr=cmove] (-4.4, 0) -- (-3.2, 0)
    node[axislabel, text=cmove, right=0.5mm] {右へ移動};

  % === Right stick: altitude + turning / 右スティック: 高度と旋回 ===
  \node[stickname] at (6.0, 4.1) {右スティック};
  \draw[cstick, line width=2pt, fill=black!25] (6.0, 0) circle (1.4);
  \draw[cstick, line width=1pt, fill=black!40] (6.0, 0) circle (0.9);

  \draw[arr=cvert] (6.0, 1.6) -- (6.0, 2.8)
    node[axislabel, text=cvert, above=0.5mm] {上昇};
  \draw[arr=cvert] (6.0, -1.6) -- (6.0, -2.8)
    node[axislabel, text=cvert, below=0.5mm] {下降};
  \draw[arr=cvert] (4.4, 0) -- (3.2, 0)
    node[axislabel, text=cvert, left=0.5mm] {左に旋回};
  \draw[arr=cvert] (7.6, 0) -- (8.8, 0)
    node[axislabel, text=cvert, right=0.5mm] {右に旋回};

  % === Notes / 補足 ===
  \node[note] at (-6.0, -4.7) {（機体の水平移動）};
  \node[note] at (6.0, -4.7) {（高度と機首の向き）\\ 押し込み = ARM／DISARM};

\end{tikzpicture}

\end{document}

  \end{adjustbox}

  \vspace{0.2em}

  {\small 今日つかう飛行モード \textbf{ALT\_HOLD}（高度維持）では高度は自動維持 --- 操縦は主に\textbf{左スティック}を少しずつ倒す}
\end{frame}

% ------------------------------------------------------------------
\begin{frame}[fragile]{① まずはシミュレータで飛ばしてみる}
  \begin{block}{コマンドプロンプト（文字で PC を操作する画面）で実行}
    \begin{lstlisting}[style=sfbash, numbers=none]
sf sim run   ← 3Dシミュレータが起動
    \end{lstlisting}
  \end{block}

  {\footnotesize ※ 「←」以降は説明なので入力しません（\code{sf} = StampFly 用コマンド）}

  \vspace{0.2em}

  \begin{columns}[c]
    \begin{column}{0.50\textwidth}
      \begin{itemize}\small
        \item コントローラを USB ケーブルで PC に接続（設定済み）
        \item 起動直後の校正が終わるまでスティックに触らない
        \item スティックの割り当ては\textbf{実機とまったく同じ}
      \end{itemize}
    \end{column}
    \begin{column}{0.48\textwidth}
      \centering
      \includegraphics[width=0.82\columnwidth]{simulator_demo.png}\\
      {\scriptsize シミュレータ画面 --- 画面の StampFly を操縦する}
    \end{column}
  \end{columns}
\end{frame}

% ------------------------------------------------------------------
\begin{frame}{① シミュレータ → 実機への切り替え}
  \begin{block}{コントローラを実機用（ESP-NOW）に戻す}
    \begin{enumerate}\small
      \item コントローラの \textbf{M5 ボタン}（画面下）でメニューを開く
      \item \textbf{「通信モード」を1回選ぶ} $\to$ USB HID から ESP-NOW（ルーターなしで機体と直接つながる無線）に切り替わり、自動で再起動する
      \item USB ケーブルを外す（実機とは電波でつながる）
    \end{enumerate}
  \end{block}

  \vspace{0.3em}

  {\footnotesize ※ 通信モードを切り替えても、ペアリングなどの設定はそのまま残ります}
\end{frame}

% ------------------------------------------------------------------
\begin{frame}{① 機体の起動と LED の見方}
  \begin{table}
    \centering\small
    \begin{tabular}{ll}
      \hline
      \textbf{LED・音} & \textbf{状態} \\
      \hline
      白色点灯 + 起動音 & 起動中 \\
      マゼンタ点滅 & ジャイロ校正中（機体に触らない） \\
      緑色常灯 + ピッ×3 & 準備完了（約10秒） \\
      \hline
    \end{tabular}
  \end{table}

  \vspace{0.5em}

  \begin{alertblock}{起動時の注意}
    平らな床に置いて電源を入れる。校正が終わる（緑常灯）まで機体を動かさない
  \end{alertblock}
\end{frame}

% ------------------------------------------------------------------
\begin{frame}{① コントローラのボタンとスティック}
  \begin{columns}[c]
    \begin{column}{0.40\textwidth}
      \centering
      \includegraphics[width=0.95\columnwidth]{atom_joystick.jpg}\\
      {\scriptsize Atom JoyStick（左上が L ボタン）}
    \end{column}
    \begin{column}{0.56\textwidth}
      \begin{table}
        \centering\footnotesize
        \begin{tabular}{ll}
          \hline
          \textbf{飛行モード} & \textbf{機体 LED 色} \\
          \hline
          STABILIZE & 黄緑 \\
          ACRO      & 青 \\
          ALT\_HOLD & オレンジ \\
          POS\_HOLD & マゼンタ \\
          \hline
        \end{tabular}
      \end{table}

      \begin{block}{ALT\_HOLD への切り替え}
        \small L ボタンを \textbf{1回} 押す（OFF $\to$ ALT\_HOLD）\\
        機体 LED が\textbf{オレンジ}になったことを確認する
      \end{block}
    \end{column}
  \end{columns}
\end{frame}

% ------------------------------------------------------------------
\begin{frame}{① ALT\_HOLD と自動離着陸}
  \begin{block}{自動離着陸と ARM / DISARM}
    \small \textbf{ARM}（右スティックを押し込む）$\to$ 自動離陸・高度0.5\,m を自動維持\\
    \textbf{DISARM}（もう一度押し込む）$\to$ 自動着陸 $\to$ 着陸中にもう一度で緊急停止
  \end{block}

  \vspace{0.3em}

  \begin{alertblock}{低電圧警告}
    \footnotesize シアン点滅 + 警告音 $\to$ 着陸して電池交換（飛行時間 約4分）
  \end{alertblock}

  \vspace{0.3em}

  \begin{exampleblock}{進め方}
    \footnotesize 1人1セット・全員同時に飛行（交代なし）。ペアリングは1対1 --- 隣の機体は反応しません。電池が切れたら（シアン点滅）交換して再開
  \end{exampleblock}
\end{frame}
